% Chapter 3: Quantum Double Schubert Polynomials
% 基于 Kirillov & Maeno (1996)
% arXiv: alg-geom/9605005

\chapter{Quantum Double Schubert Polynomials}
\label{chap:QuantumDouble}

\section{引言：从经典几何到曲线计数——为什么要引入量子版本？}
\label{sec:QuantumMotivation}

\subsection{经典相交理论的局限}

在前两章中，我们研究了经典双重 Schubert 多项式 $\mathfrak{S}_w(\mathbf{x}; \mathbf{y})$ 的性质，以及三重变量集情形下的 Graham positivity。

但当我们研究旗流形（flag manifold）上曲线的个数时，普通的相交理论遇到了根本性的困难。

设想我们想要计数 $\mathbb{P}^1$ 到旗流形的稳定映射个数。在经典相交理论中，我们只能计数点与点的交集——这是一种静态的计数。但带参数的曲线涉及额外的自由度：每条曲线有自己的"度数"（degree），这对应于同调类 $\beta \in H_2(Fl_n, \mathbb{Z})$。

这引导我们进入\textbf{量子上同调}（quantum cohomology）的世界——一个源于物理中拓扑场论（topological field theory）的思想，由 Kontsevich 和 Manin 系统化。\footnote{问：什么是量子上同调？它与经典上同调的本质区别是什么？为什么量子版本来自物理学？见附录 \cref{sec:quantum-cohomology-origin}。}

\subsection{量子上同调的核心思想}

在量子上同调（quantum cohomology）中，Schubert 类的乘法不再是普通的乘积，而是一种带有形式变量 $q^\beta$ 的新运算
\[
[\sigma_u] \star [\sigma_v] = \sum_{w,\beta} q^{\beta} c^w_{u,v}(\beta) [\sigma_w].
\label{eq:QuantumSchubertMultiplication}
\]
其中 $\beta$ 记录了连接子簇的曲线类（curve class），$q^{\beta}$ 是形式变量。这一公式最初来自物理中的 \textbf{A-twisted topological field theory}，它揭示了曲线计数与多项式结构常数之间的深层联系。\footnote{公式 \eqref{eq:QuantumSchubertMultiplication} 的详细解释和几何意义见附录 \cref{sec:quantum-multiplication-meaning}。}

\textbf{关键洞察}：量子上同调中的乘法 $q^\beta$ 因子编码了曲线计数的额外信息——这正是经典情形所缺失的。

Kirillov 和 Maeno 在 1996 年的工作中引入了\textbf{量子双重 Schubert 多项式} $\widetilde{\mathfrak{S}}_w(\mathbf{x}; \mathbf{y}; \mathbf{q})$，这是旗流形等变量量子上同调类的 Lascoux-Schützenberger 型代表元。本文介绍这一工作的核心内容。\footnote{Kirillov-Maeno 1996 年的工作背景及其与先前工作的关系见附录 \cref{sec:KM1996-background}。}

\section{经典回顾：除差算子与 Schubert 多项式}
\label{sec:ClassicalRecap}

在引入量子版本之前，我们先简要回顾经典理论的核心定义——这些内容在 Chapter 1 中已有详细讨论，这里做快速梳理以保持连贯性。

\subsection{除差算子}

除差算子（divided difference operator）是 Schubert 演算的核心工具。我们在 Chapter 1 中已经看到，$\partial_i$ 测量了多项式在交换两个变量时的变化率：

\begin{Definition}[除差算子]
\label{def:DividedDifferenceOp}
对于 $i = 1, \ldots, n-1$，\textbf{除差算子} $\partial_i$ 定义为：
\begin{equation}
\partial_i f = \frac{f(x_1,\ldots,x_i,x_{i+1},\ldots) - f(x_1,\ldots,x_{i+1},x_i,\ldots)}{x_i - x_{i+1}}.
\label{eq:DividedDifference}
\end{equation}
\end{Definition}

\textbf{关键性质}：$\partial_i$ 满足 Leibniz 法则
\begin{equation}
\partial_i(fg) = (\partial_i f)g + (s_i f)(\partial_i g),
\label{eq:LeibnizRule}
\end{equation}
其中 $s_i f = f(x_1,\ldots,x_{i+1},x_i,\ldots)$ 表示交换 $x_i$ 和 $x_{i+1}$。

\subsection{经典 Schubert 多项式}

正如 Chapter 1 中所定义，Schubert 多项式通过除差算子作用在特殊多项式 $x^\delta$ 上得到：

\begin{Definition}[Schubert 多项式]
\label{def:ClassicalSchubert}
设 $\delta = (n-1,n-2,\ldots,1,0)$，$w_0$ 为 $S_n$ 的最长置换。\textbf{Schubert 多项式}定义为：
\begin{equation}
\mathfrak{S}_w(x) = \partial_{w^{-1}w_0}(x^{\delta}),
\label{eq:SchubertPolynomial}
\end{equation}
其中 $x^{\delta} = x_1^{n-1}x_2^{n-2}\cdots x_n^0$。
\end{Definition}

当引入 $y$ 变量后，我们得到双重版本——这正是 Chapter 2 研究乘积公式时使用的工具：

\begin{Definition}[双重 Schubert 多项式]
\label{def:DoubleSchubert}
\textbf{双重 Schubert 多项式}定义为：
\begin{equation}
\mathfrak{S}_w(x,y) = \partial_{w^{-1}w_0}^{(x)} \prod_{i+j\leq n}(x_i + y_j).
\label{eq:DoubleSchubert}
\end{equation}
\end{Definition}

\section{量子上同调的基本框架}
\label{sec:QuantumCohomologyFramework}

为什么量子上同调是自然的？这需要从几何学中寻求解释。

在经典相交理论中，我们计算子流形 $X, Y \subset V$ 的交集 $X \cap Y$ 的点数。但在旗流形的情形，当我们要计数\textbf{连接两点的曲线}时，普通的相交理论失效了——因为这不是点与点的交集，而是曲线与曲线的交汇。

Kontsevich 和 Manin 的伟大贡献在于，他们意识到这种计数可以系统化——通过引入 Gromov-Witten 不变量（Gromov-Witten invariant）。\footnote{Gromov-Witten 不变量的详细定义和计算方法见附录 \cref{sec:GW-invariant-detailed}。}

\subsection{Gromov-Witten 不变量（Gromov-Witten invariant）}

设 $V$ 是光滑复流形，$\beta \in H_2(V, \mathbb{Z})$ 是曲线类（curve class）。\textbf{Gromov-Witten 不变量}
\[
\langle I_{0,m,\beta}^V \rangle : H^*(V, \mathbf{Q})^{\otimes m} \longrightarrow \mathbf{Q}.
\label{eq:GromovWittenInvariant}
\]
计数满足特定条件的参数化曲线个数。

\textbf{直观理解}：$ \langle I_{0,m,\beta}^V \rangle(X_1^* \otimes \cdots \otimes X_m^*) $ 是满足以下条件的稳定映射（stable map） $f: \mathbb{P}^1 \to V$ 的"虚拟个数"：
\begin{itemize}
\item $f$ 的像代表同调类 $\beta$；
\item $f(p_i) \in X_i$（其中 $p_1, \ldots, p_m \in \mathbb{P}^1$ 是标记点）。
\end{itemize}

这些不变量最初来自物理中的拓扑场论，后来被严格化——这又是数学从物理学汲取营养的又一经典案例。\footnote{问：Gromov-Witten 不变量最初是如何从物理中推导出来的？A-model 和 B-model 是什么？见附录 \cref{sec:GW-physics-origin}。}

\subsection{量子上同调的几何意义}

理解量子上同调的关键在于认识到它编码了\textbf{曲线计数}的额外信息。

在经典 Schubert calculus 中，我们问的是："两个 Schubert 簇 $X(u)$ 和 $X(v)$ 在旗流形中相交于几个点？"

在量子上同调中，我们问的是更深的问题："给定同调类 $\beta$，有多少条曲线连接 $X(u)$ 和 $X(v)$，并且这些曲线通过了旗流形中的特定位置？"

这类似于从问"两点之间有多少条直线"到问"两点之间有多少条抛物线"——问题变得更加丰富和复杂。

\textbf{为什么这是重要的？} 因为：
\begin{enumerate}
\item 量子上同调包含了经典理论无法触及的代数几何结构；
\item 量子修正项揭示了旗流形的更深层性质；
\item 量子上同调与弦理论的对偶性提供了新的计算工具。
\end{enumerate}

\textbf{量子变量的含义}：形式变量 $q_i$ 编码了第 $i$ 个曲线类的信息。当我们说某个系数"包含 $q_1^2 q_3$ 项"时，意思是这个项对应着两条 1 类曲线和一条 3 类曲线的组合。

\begin{Example}[$\mathbb{P}^1$ 上的曲线计数]
最简单的例子是 $\mathbb{P}^1$（射影直线）。在经典相交理论中：
\[
[\text{点}] \cdot [\text{点}] = 1 \cdot [\mathbb{P}^1]
\]
因为两个点确定一条直线。

在量子上同调中，我们需要考虑曲线的亏格（genus）和次数。但对于 $\mathbb{P}^1$（亏格 0），量子修正是平凡的——这体现了量子效应的微妙之处。
\end{Example}

\begin{Example}[旗流形上的曲线计数]
在旗流形 $Fl_n$ 中，计数曲线是一个深刻的问题。不同的同调类 $\beta$ 对应于旗流形中不同"方向"的曲线族。

例如，在 $Fl_3$ 中，从某个旗到另一个旗的曲线可能属于不同的同调类。量子上同调中的 $q$ 变量正是记录了这些不同曲线类的相对权重。
\end{Example}

\subsection{小量子上同调环（small quantum cohomology ring）}

对于旗流形（flag manifold） $Fl_n = SL_n / B$，小量子上同调环（small quantum cohomology ring）的结构由 Givental、Kim 以及 Ciocan-Fontanine 确定。\footnote{各数学家贡献的历史顺序和相互关系见附录 \cref{sec:quantum-cohomology-history}。}

\begin{Theorem}[小量子上同调环 {\cite[ Theorem 8]{KM}}]
\label{th:SmallQuantumCohomology}
小量子上同调环 $QH^*(Fl_n)$ 同构于
\[
\mathbf{Z}[x_1, \ldots, x_n; q_1, \ldots, q_{n-1}] / (\widetilde{e}_1, \ldots, \widetilde{e}_n).
\label{eq:SmallQuantumCohomologyRing}
\]
其中 $\widetilde{e}_i(x|q)$ 是\textbf{量子初等对称函数}，由以下行列式生成：
\[
\det\begin{pmatrix}
x_1 + t & q_1 & 0 & \cdots & 0 \\
-1 & x_2 + t & q_2 & \cdots & 0 \\
0 & -1 & x_3 + t & \cdots & 0 \\
\vdots & \ddots & \ddots & \ddots & \vdots \\
0 & \cdots & 0 & -1 & x_n + t
\end{pmatrix}
= t^n + \widetilde{e}_1 t^{n-1} + \cdots + \widetilde{e}_n.
\label{eq:QuantumElementaryGeneratingFunction}
\]
\end{Theorem}

当 $q_1 = \cdots = q_{n-1} = 0$ 时，\eqref{eq:SmallQuantumCohomologyRing} 退化为经典上同调环——这体现了量子版本与经典版本之间的深刻联系。

\textbf{几何含义}：量子初等对称函数 $\widetilde{e}_i$ 对应于旗流形中特定子流形的类。它们的零点描述了旗流形的量子修正结构。

\section{量子双重 Schubert 多项式的定义}
\label{sec:DefinitionQuantumDouble}

有了量子上同调框架后，我们现在可以正式定义量子双重 Schubert 多项式。

\subsection{顶部多项式（top polynomial）}

设 $x = (x_1, \ldots, x_n)$ 和 $y = (y_1, \ldots, y_n)$ 是两组变量，$q = (q_1, \ldots, q_{n-1})$ 是量子变量。

\textbf{顶部多项式}（top polynomial）定义为
\[
\widetilde{\mathfrak{S}}_{w_0}(x,y) := \prod_{i=1}^{n-1} \Delta_i(y_{n-i} \mid x_1, \ldots, x_i).
\label{eq:TopPolynomial}
\]
其中
\[
\Delta_k(t \mid x_1, \ldots, x_k) := \sum_{j=0}^k t^{k-j} e_j(x_1, \ldots, x_k \mid q_1, \ldots, q_{k-1}).
\label{eq:DeltaFunction}
\]
是关于 $x_1, \ldots, x_k$ 的量子初等对称函数的生成函数。\footnote{量子初等对称函数 $e_j(x_1, \ldots, x_k \mid q_1, \ldots, q_{k-1})$ 的显式公式和递推关系见附录 \cref{sec:quantum-elementary-explicit}。}

\textbf{为什么叫"顶部多项式"？} 当 $q=0$ 时，$\widetilde{\mathfrak{S}}_{w_0}(x,y) = \prod_{i+j\leq n}(x_i+y_j)$，这正是对应最长置换 $w_0$ 的 Schubert 多项式——在 Bruhat 序（Bruhat order）中，它是旗流形的"最大维" Schubert 胞腔。

\subsection{量子双重 Schubert 多项式}

有了顶部多项式，我们通过对它施加除差算子来生成所有量子双重 Schubert 多项式：

\begin{Definition}[量子双重 Schubert 多项式 {\cite[ Definition 4]{KM}}]
\label{def:QuantumDoubleSchubert}
对于每个置换 $w \in S_n$，定义\textbf{量子双重 Schubert 多项式}为
\[
\widetilde{\mathfrak{S}}_w(x,y) = \partial_{ww_0}^{(y)} \widetilde{\mathfrak{S}}_{w_0}(x,y).
\label{eq:QuantumDoubleSchubertDef}
\]
其中除差算子 $\partial_{ww_0}^{(y)}$ 作用于 $y$ 变量。
\end{Definition}

\textbf{与经典版本的关系}：当 $q = 0$ 时，
\[
\left. \widetilde{\mathfrak{S}}_w(x,y) \right|_{q=0} = \mathfrak{S}_w(x,y).
\label{eq:ClassicalLimit}
\]
即经典双重 Schubert 多项式。这是"经典极限"的一个典型例子——量子理论在 $q \to 0$ 时退化为经典理论。

验证：设 $q=0$，则 $\Delta_i(y_{n-i} \mid x_1, \ldots, x_i) = \prod_{j=1}^i (x_j + y_{n-i+j})$，而
\[
\prod_{i=1}^{n-1} \prod_{j=1}^i (x_j + y_{n-i+j}) = \prod_{i+j \leq n}(x_i + y_j) = \mathfrak{S}_{w_0}(x,y).
\]
然后由除差算子的性质可得 \eqref{eq:QuantumDoubleSchubertDef}。

\section{主要性质}
\label{sec:MainProperties}

量子双重 Schubert 多项式具有一系列重要性质。

\subsection{稳定性（stability）}

与经典版本类似（见 Chapter 1 的稳定性（stability）讨论），量子双重 Schubert 多项式也具有稳定性：

\begin{Proposition}[稳定性]
\label{prop:QuantumStability}
设 $m > n$ 且 $i: S_n \hookrightarrow S_m$ 是自然嵌入，则
\[
\widetilde{\mathfrak{S}}_w(x,y) = \widetilde{\mathfrak{S}}_{i(w)}(x,y).
\label{eq:QuantumStability}
\]
\end{Proposition}

这意味着量子双重 Schubert 多项式在嵌入下保持不变——我们可以将其定义在 $S_\infty$ 上，这对于理论研究非常便利。

\textbf{为什么稳定性重要？} 在经典 Schubert 多项式理论中，稳定性保证了多项式可以作为 $S_\infty$ 上的对象来研究，而不必关心具体的 $n$。量子版本继承了这一性质，使得无穷维的分析成为可能。

\subsection{$S_3$ 的完整例子}

为了建立直观感受，我们展示 $S_3$ 的量子双重 Schubert 多项式的完整集合：

\begin{Example}[$S_3$ 的量子双重 Schubert 多项式]
在 $S_3$ 中，量子双重 Schubert 多项式为：
\begin{align*}
\widetilde{\mathfrak{S}}_{s_1s_2s_1}(\mathbf{x}; \mathbf{y}) &= (x_1+y_2)(x_1+y_1)(x_2+y_1) + q_1(x_1+y_2), \\
\widetilde{\mathfrak{S}}_{s_2s_1}(\mathbf{x}; \mathbf{y}) &= (x_1+y_1)(x_1+y_2) - q_1, \\
\widetilde{\mathfrak{S}}_{s_1s_2}(\mathbf{x}; \mathbf{y}) &= (x_1+y_1)(x_2+y_1) + q_1, \\
\widetilde{\mathfrak{S}}_{s_1}(\mathbf{x}; \mathbf{y}) &= x_1 + y_1, \\
\widetilde{\mathfrak{S}}_{s_2}(\mathbf{x}; \mathbf{y}) &= x_1 + x_2 + y_1 + y_2, \\
\widetilde{\mathfrak{S}}_{\mathrm{id}}(\mathbf{x}; \mathbf{y}) &= 1.
\end{align*}
\end{Example}

\textbf{观察}：每个多项式都由一个经典部分加上一到两个 $q$ 修正项组成。当 $q_1 = 0$ 时，这些公式退化为经典双重 Schubert 多项式。

\textbf{关键洞察}：第一个公式（对应最长元 $w_0 = s_1s_2s_1$）包含了 $q_1$ 项，这反映了在量子上同调中，最高维 Schubert 类的乘法会涉及量子修正——这正是量子与经典理论的核心差异所在。

\begin{Example}[验证经典极限]
以 $\widetilde{\mathfrak{S}}_{s_2s_1}$ 为例。当 $q_1 = 0$ 时，
$$\widetilde{\mathfrak{S}}_{s_2s_1} = (x_1+y_1)(x_1+y_2)$$

由第一章的知识，经典双重 Schubert 多项式满足
$$\mathfrak{S}_{s_2s_1}(x,y) = (x_1+y_1)(x_1+y_2) = x_1^2 + x_1y_1 + x_1y_2 + y_1y_2$$

这正是 $(x_1+y_1)(x_1+y_2)$ 的展开。经典极限验证通过。
\end{Example}

\section{量子 Schubert 多项式}
\label{sec:QuantumSchubertPolynomials}

除了量子双重版本，Kirillov 和 Maeno 还定义了不含 $y$ 变量的量子 Schubert 多项式。

\subsection{量子配对（quantum pairing）}

在经典理论中，我们用标量积 $\langle f, g \rangle = \partial_{w_0}(fg)$ 来定义 Schubert 多项式的正交性（orthogonality）。在量子情形，我们需要用量子配对（quantum pairing）来替代：

定义量子配对（quantum pairing）为 Grothendieck 留数（Grothendieck residue）：
\[
\langle f, g \rangle_Q = \mathrm{Res}_{\widetilde{I}}(fg).
\label{eq:QuantumPairing}
\]
其中 $\widetilde{I}$ 是由量子初等对称函数生成的理想。这个配对对应于量子上同调环中的相交形式（intersection form）。

\textbf{为什么使用 Grothendieck 留数（Grothendieck residue）？} 这是因为在商环 $QH^*(Fl_n)$ 中，普通的积分被留数公式所替代——这反映了量子版本与经典版本的本质差异。\footnote{Grothendieck 留数的定义和性质见附录 \cref{sec:grothendieck-residue}。}

\subsection{定义}

\begin{Definition}[量子 Schubert 多项式 {\cite[ Definition 5]{KM}}]
\label{def:QuantumSchubert}
\textbf{量子 Schubert 多项式} $\widetilde{\mathfrak{S}}_w(x)$ 是将字典序排列的单项式 $\{x^I \mid I \subset \delta\}$ 关于量子配对 $\langle, \rangle_Q$ 进行 Gram-Schmidt 正交化（Gram-Schmidt orthogonalization）得到的。
\end{Definition}

这一定义在结构上与经典 Schubert 多项式的刻画像（见 Chapter 2 的讨论）完全对应——只是将普通标量积替换为了量子标量积。

量子 Schubert 多项式满足以下正交性（orthogonality）条件：
\[
\langle \widetilde{\mathfrak{S}}_u, \widetilde{\mathfrak{S}}_v \rangle_Q = \begin{cases}
1 & \text{若 } v = w_0 u, \\
0 & \text{其他}.
\end{cases}
\label{eq:QuantumSchubertOrthogonality}
\]

\textbf{几何意义}：这一正交性对应于量子上同调环中 Schubert 类的相交数。在量子版本中，相交不再只是点计数，而是涉及曲线的虚拟计数。

\section{量子 Cauchy 恒等式}
\label{sec:QuantumCauchy}

这是本文的核心结果之一——它建立了量子与经典 Schubert 多项式之间的对偶关系。

\begin{Theorem}[量子 Cauchy 恒等式 {\cite[ Theorem B]{KM}}]
\label{th:QuantumCauchy}
设 $\widetilde{\mathfrak{S}}_w(x) = \widetilde{\mathfrak{S}}_w(x,0)$，则
\[
\sum_{w \in S_n} \widetilde{\mathfrak{S}}_w(x) \mathfrak{S}_{ww_0}(y) = \widetilde{\mathfrak{S}}_{w_0}(x,y).
\label{eq:QuantumCauchy}
\]
\end{Theorem}

\textbf{这为什么重要？} 这一恒等式将量子 Schubert 多项式与经典 Schubert 多项式联系起来。当 $q = 0$ 时，\eqref{eq:QuantumCauchy} 退化为经典的 Cauchy 公式
\[
\sum_{w \in S_n} \mathfrak{S}_w(x) \mathfrak{S}_{ww_0}(y) = \prod_{i+j \leq n}(x_i + y_j).
\label{eq:ClassicalCauchy}
\]
这正是 Chapter 2 中研究乘积公式时的基础。

\textbf{几何意义}：\eqref{eq:QuantumCauchy} 建立了量子 Schubert 多项式（对应量子上同调类）与经典 Schubert 多项式（对应普通上同调类）之间的对偶关系。这种对偶性在数学物理中有深刻的根源。\footnote{量子 Cauchy 恒等式的物理起源（拓扑场论中的对偶性）见附录 \cref{sec:cauchy-physics-origin}。}

\section{Lascoux-Schützenberger 型公式}
\label{sec:LascouxFormula}

作为量子 Cauchy 恒等式的一个直接推论，我们得到量子 Schubert 多项式的 Lascoux-Schützenberger 型公式。

\begin{Corollary}[Lascoux-Schützenberger 型公式 {\cite[ Theorem C]{KM}}]
\label{cor:QuantumLascoux}
设 $\widetilde{\mathfrak{S}}_{w_0}(x,y)$ 如上定义，则
\[
\widetilde{\mathfrak{S}}_w(x) = \partial_{ww_0}^{(y)} \widetilde{\mathfrak{S}}_{w_0}(x,y) \big|_{y=0}.
\label{eq:QuantumLascouxFormula}
\]
\end{Corollary}

\textbf{这告诉我们什么？} 量子 Schubert 多项式可以通过对量子双重 Schubert 多项式施加除差算子，然后取 $y=0$ 来得到。这与经典情形
\[
\mathfrak{S}_w(x) = \partial_{w^{-1}w_0} \mathfrak{S}_{w_0}(x).
\label{eq:ClassicalLascouxFormula}
\]
的结构完全类似——量子版本只是经典版本的自然推广。

\section{与三重 Schubert Positivity 的联系}
\label{sec:ConnectionTriplePositivity}

在 Chapter 1 中，我们研究了 Gao-Xiong 关于三重 Schubert Positivity（triple Schubert positivity）的结果。现在我们揭示量子双重 Schubert 多项式与这一结果的联系。

回顾 Chapter 1 的核心问题：研究展开系数
\[
\mathfrak{S}_u(\mathbf{x}; \mathbf{y}) \cdot \mathfrak{S}_v(\mathbf{x}; \mathbf{t}) = \sum_w c^w_{u,v}(\mathbf{y}, \mathbf{t}) \mathfrak{S}_w(\mathbf{x}; \mathbf{t}).
\label{eq:TripleProductExpansion}
\]
的正性。

\textbf{关键联系}：当我们考虑量子上同调时，同样的乘法在量子配对下有更丰富的结构。量子双重 Schubert 多项式提供了这些乘法结构的代数代表元。

具体而言，等变量量子上同调（equivariant quantum cohomology）环 $QH^*_{U_n}(Fl_n)$ 同构于
\[
\mathbf{Z}[x_1, \ldots, x_n; y_1, \ldots, y_n; q_1, \ldots, q_{n-1}] / \widetilde{J}.
\label{eq:EquivariantQuantumCohomology}
\]
其中理想 $\widetilde{J}$ 由
\[
e_i(x \mid q_1, \ldots, q_{n-1}) - e_i(y_1, \ldots, y_n), \quad 1 \leq i \leq n.
\label{eq:QuantumIdealRelations}
\]
生成。

量子双重 Schubert 多项式 $\widetilde{\mathfrak{S}}_w(x,y)$ 在这个环中扮演的角色，与经典双重 Schubert 多项式在经典等变量子上同调中扮演的角色完全类似。

\textbf{三重版本的关键洞察}：当 Gao-Xiong 研究三重变量集 $(x, y, t)$ 时，他们实际上是在研究等变量量子上同调（equivariant quantum cohomology）的一个精化版本。量子双重 Schubert 多项式的结构常数提供了这些精化的代数解释。

\section{Vafa-Intriligator 公式}
\label{sec:VafaIntriligator}

本文的另一重要贡献是给出了 Vafa-Intriligator 公式（Vafa-Intriligator formula）的高维推广。

Vafa-Intriligator 公式最初来自物理，它给出了某些余维一子流形上交点数的显式计算。在旗流形的情形，这涉及 Gromov-Witten 不变量的显式公式。

Kirillov 和 Maeno 在 Section 8 中给出了旗流形的 Vafa-Intriligator 公式的高维推广。这一公式在量子上同调中的地位，如同经典 Schubert 多项式在普通上同调中的地位——它提供了计算结构常数的组合工具。

我们将在附录中给出这一公式的详细推导（见 \cref{sec:appendix-ch3-vafa-intriligator}）。

\section{量子化映射}
\label{sec:QuantizationMap}

Kirillov 和 Maeno 还引入了量子化映射，这是理解量子与经典关系的关键工具。

定义量子化映射 $P_n \to \overline{P}_n$ 为
\[
\widetilde{f}(x) = \sum_{w \in S^{(n)}} \partial_w^{(y)} f(y) \widetilde{\mathfrak{S}}_w(x,y) \big|_{\overline{P}_n}.
\label{eq:QuantizationMap}
\]

\begin{Proposition}[配对保持（pairing preservation）]
\label{prop:PairingPreservation}
量子化映射保持配对（pairing preservation）：
\[
\langle \widetilde{f}, \widetilde{g} \rangle_Q = \langle f, g \rangle, \quad f, g \in P_n.
\label{eq:PairingPreservation}
\]
\end{Proposition}

这意味着量子化是配对保持的线性映射，但\textbf{不保持乘法}（见 \cite[ Remark 5(i)]{KM}）。

量子化保持初等多项式和完全齐次对称多项式，但将单形式映射到量子单形式。

\textbf{为什么这个性质重要？} 配对保持意味着量子化保持了相交理论的核心结构——但乘法结构的变化正是量子效应的体现。这揭示了量子与经典之间的深刻联系和本质差异。

\section{与其他工作的关系}
\label{sec:RelationOtherWork}

本文与 Fomin、Gelfand 和 Postnikov 的独立工作有部分重叠，但采用了完全不同的方法。

Postnikov 等人基于一类交换算子 $X_i$ 发展了量子 Schubert 多项式理论，得到了量子 Monk 公式、正交性和量子乘法的定义等结果。

Kirillov-Maeno 的方法则基于量子 Cauchy 恒等式和留数配对，提供了另一种理解量子 Schubert 结构常数的途径。这两种方法互为补充，共同构成了量子 Schubert calculus（quantum Schubert calculus）的基础。\footnote{Fomin-Gelfand-Postnikov 方法与 Kirillov-Maeno 方法的比较见附录 \cref{sec:KM-vs-FGP}。}

\section{本章小结}
\label{sec:Chapter3Summary}

本章介绍了 Kirillov 和 Maeno 关于量子双重 Schubert 多项式的工作。关键要点如下：

\begin{enumerate}
\item \textbf{量子上同调的动机}：经典相交理论无法处理曲线计数问题，需要引入量子修正
\item \textbf{量子初等对称函数} $\widetilde{e}_i(x|q)$ 由行列式生成函数定义，它们生成了小量子上同调环的 Relations
\item \textbf{量子双重 Schubert 多项式} $\widetilde{\mathfrak{S}}_w(x,y)$ 通过对顶部多项式施加除差算子定义
\item \textbf{经典极限}：当 $q=0$ 时，$\widetilde{\mathfrak{S}}_w(x,y)$ 退化为经典双重 Schubert 多项式
\item \textbf{量子 Cauchy 恒等式}建立了量子 Schubert 多项式与经典 Schubert 多项式之间的对偶关系
\item \textbf{量子化映射}是配对保持的线性映射，但不保持乘法
\item 这些结果与 Chapter 1 的 Gao-Xiong 三重 Schubert positivity 工作有深刻的联系
\end{enumerate}

在下一章中，我们将继续探讨 Mihalcea 关于等变量子量子上同调正性的结果。

\section{附录：量子初等对称函数的显式公式}
\label{sec:quantum-elementary-explicit}

\subsection*{背景}
量子初等对称函数 $\widetilde{e}_i(x|q)$ 通过行列式生成函数定义。本节给出其显式展开形式。

\subsection*{行列式递推}

考虑 $n \times n$ 三对角矩阵
$$A = \begin{pmatrix}
x_1 + t & q_1 & 0 & \cdots & 0 \\
-1 & x_2 + t & q_2 & \cdots & 0 \\
0 & -1 & x_3 + t & \cdots & 0 \\
\vdots & \ddots & \ddots & \ddots & \vdots \\
0 & \cdots & 0 & -1 & x_n + t
\end{pmatrix}$$

定义 $D_m(t) = \det(A_m)$，其中 $A_m$ 是 $m \times m$ 的主子矩阵。则 $D_m(t)$ 满足递推关系：
$$D_m(t) = (x_m + t) D_{m-1}(t) - q_{m-1} D_{m-2}(t)$$
其中 $D_0(t) = 1$，$D_1(t) = x_1 + t$。

\subsection*{显式公式}

通过递推可得：
$$\widetilde{e}_k(x|q) = \sum_{j=0}^{\lfloor k/2 \rfloor} (-1)^j \sum_{1 \leq i_1 < \cdots < i_j \leq k-1} q_{i_1} q_{i_2} \cdots q_{i_j} \cdot e_{k-2j}(x_1, \ldots, x_{k-j}).$$

当 $q=0$ 时，这简化为 $e_k(x_1, \ldots, x_k)$。

\section{附录：Vafa-Intriligator 公式的详细推导}
\label{sec:appendix-ch3-vafa-intriligator}

\subsection*{历史背景}

Vafa 和 Intriligator 在 1996 年基于物理中的拓扑场论提出了一个计算某些余维一子流形上交点数目的公式。

在旗流形的情形，这个公式涉及 Gromov-Witten 不变量的显式计算。

\subsection*{公式陈述}

对于 Grassmannian $G/P$，Vafa-Intriligator 公式给出：
$$c_{u,v}^w(q) = \sum_{\beta} q^\beta \langle I_{0,3,\beta}(\sigma_u, \sigma_v, \sigma_{w^\vee}) \rangle$$

其中求和遍历所有有效的曲线类 $\beta$。

详细证明涉及以下步骤（详见 \cite[ Section 8]{KM}）：
\begin{enumerate}
\item 将问题化为留数计算
\item 应用 Localization 定理
\item 验证与已知结果的一致性
\end{enumerate}

\section{附录：Grothendieck 留数的定义}
\label{sec:grothendieck-residue}

\subsection*{定义}

设 $\widetilde{I}_n \subset \overline{P}_n = \mathbf{Z}[q_1,\ldots,q_{n-1}][x_1,\ldots,x_n]$ 是由量子初等对称函数生成的理想。设 $\mathcal{H} = \det\left(\frac{\partial \widetilde{e}_i}{\partial x_j}\right)$ 是 Hessian。

对于 $f \in \overline{P}_n$：
\begin{itemize}
\item 若 $\deg f < \deg \mathcal{H}$，则 $\mathrm{Res}_{\widetilde{I}}(f) = 0$
\item 若 $\deg f = \deg \mathcal{H}$，则 $f \equiv \frac{\alpha}{n!}\mathcal{H} \pmod{\widetilde{I}_n}$，其中 $\mathrm{Res}_{\widetilde{I}}(f) = \alpha$
\end{itemize}

\subsection*{在量子配对中的应用}

配对定义为 $\langle f, g \rangle_Q = \mathrm{Res}_{\widetilde{I}}(fg)$。这个配对在商环 $\overline{P}_n / \widetilde{I}_n \cong QH^*(Fl_n)$ 上诱导非退化配对。

\section{附录：量子 Cauchy 恒等式的证明思路}
\label{sec:cauchy-proof-sketch}

量子 Cauchy 恒等式
$$\sum_{w \in S_n} \widetilde{\mathfrak{S}}_w(x) \mathfrak{S}_{ww_0}(y) = \widetilde{\mathfrak{S}}_{w_0}(x,y)$$
的证明分为以下步骤：

步骤 1：两边同时关于 $x$ 变量应用除差算子 $\partial_{vw_0}^{(x)}$。

步骤 2：利用正交性关系简化左边。

步骤 3：验证所得表达式等于右边。

详细证明见 \cite[ Theorem B]{KM}。

\section{附录：两种量子化方法的比较}
\label{sec:KM-vs-FGP}

Kirillov-Maeno 方法和 Fomin-Gelfand-Postnikov 方法是量子 Schubert 演算中两条独立发展的路线。

\textbf{Kirillov-Maeno 方法}：
\begin{itemize}
\item 核心工具：量子 Cauchy 恒等式和 Grothendieck 留数配对
\item 多项式对象：量子双重 Schubert 多项式 $\widetilde{\mathfrak{S}}_w(x,y)$
\item 优点：与经典理论的结构对应清晰
\end{itemize}

\textbf{Fomin-Gelfand-Postnikov 方法}：
\begin{itemize}
\item 核心工具：交换算子 $X_i$
\item 多项式对象：量子 Schubert 多项式 $Q_w(x;q)$
\item 优点：组合结构更显式
\end{itemize}

两种方法在 $q=0$ 时都退化为经典 Schubert 理论。

\section{附录：量子上同调与经典上同调的比较表}
\label{sec:quantum-vs-classical-table}

\begin{center}
\begin{tabular}{c|c|c}
性质 & 经典上同调 & 量子上同调 \\ \hline
环结构 & $H^*(Fl_n)$ & $QH^*(Fl_n)$ \\
乘法 & Schubert 类相交 & 曲线计数修正 \\
参数 & 无 & 形式变量 $q$ \\
正性 & 整数系数 & 多项式系数（有待第四章）\\
配对 & 积分 & Grothendieck 留数
\end{tabular}
\end{center}

\section{附录：Gromov-Witten 不变量的详细定义}
\label{sec:GW-invariant-detailed}

\subsection*{背景}

Gromov-Witten 不变量起源于 1990 年代代数几何与拓扑弦理论的对偶研究。Gromov（1985）首先提出了紧致化轨迹空间的想法，后来被 Kontsevich（1994）系统化。

\subsection*{稳定映射空间}

设 $V$ 是光滑复流形，$\beta \in H_2(V, \mathbb{Z})$ 是曲线类。\textbf{稳定映射} $f: \mathbb{P}^1 \to V$ 是满足以下条件的映射：
\begin{itemize}
\item $f_*[\mathbb{P}^1] = \beta$（映射代表同调类 $\beta$）
\item $f$ 在标记点 $p_1, \ldots, p_m$ 处有良好定义
\item 稳定性条件：有限自同构群
\end{itemize}

稳定映射的模空间记作 $\overline{M}_{0,m}(V, \beta)$，这是一个紧致化的轨迹空间。

\subsection*{不变量的定义}

\textbf{Gromov-Witten 不变量}定义为：
$$\langle I_{0,m,\beta}^V \rangle(\gamma_1 \otimes \cdots \otimes \gamma_m) = \int_{[\overline{M}_{0,m}(V,\beta)]^{\text{vir}}} \prod_{i=1}^m \text{ev}_i^*(\gamma_i)$$

其中：
\begin{itemize}
\item $[\overline{M}_{0,m}(V,\beta)]^{\text{vir}}$ 是虚拟基本类（virtual fundamental class）
\item $\text{ev}_i: \overline{M}_{0,m}(V,\beta) \to V$ 是评估映射
\item $\gamma_i \in H^*(V, \mathbb{Q})$ 是上同调类
\end{itemize}

\subsection*{与曲线计数的关系}

当所有类都是余维一子流形的基本类时，Gromov-Witten 不变量计数的是满足条件的稳定映射的"虚拟个数"。这推广了经典的相交理论。

\section{附录：量子上同调的物理起源}
\label{sec:quantum-cohomology-origin}

\subsection*{拓扑场论背景}

量子上同调起源于物理中的\textbf{拓扑 A-模型}（topological A-model）。在弦论中，我们考虑将二维世界面（worldsheet）映射到目标空间 $V$。

当世界面是球面（亏格 0）时，对应的就是 \textbf{量子上同调}中的计数问题。

\subsection*{A-twisted topological field theory}

Witten 首先观察到，通过对 N=2 超对称场论进行 A-twist，可以得到一个拓扑场论，其算子代数正好是量子上同调环。

物理中的关联函数（correlation function）对应于数学中的 Gromov-Witten 不变量：
$$\langle \phi_1 \phi_2 \cdots \phi_m \rangle = \int_{\overline{M}_{0,m}} \prod_i \text{ev}_i^*(\phi_i)$$

\subsection*{为什么来自物理？}

数学家最初研究的是经典相交理论。但物理学家在研究拓扑弦论时发现，仅仅计数点是不够的——必须考虑世界面上的动力学。这导致了量子上同调的产生。

这是一个数学从物理学汲取深刻思想的经典案例。

\section{附录：量子多重性的几何意义}
\label{sec:quantum-multiplicity-meaning}

在公式
$$[\sigma_u] \star [\sigma_v] = \sum_{w,\beta} q^{\beta} c^w_{u,v}(\beta) [\sigma_w]$$
中，$\beta$ 编码了曲线类的信息。

\textbf{几何解释}：考虑旗流形 $Fl_n$ 中三个 Schubert 子簇 $X_u, X_v, X_w$。经典情形下，计数的是点与点的交集。但量子情形下，我们计数的是连接这些子簇的曲线。

具体而言，$c^w_{u,v}(\beta)$ 计数的是：
\begin{itemize}
\item 从 $X_u$ 中一点出发
\item 到 $X_v$ 中一点
\item 类为 $\beta$ 的曲线
\item 终点在 $X_w$ 中
\end{itemize}

当 $\beta = 0$（即点映射）时，这退化为经典相交数。

\section{附录：旗流形的 Vafa-Intriligator 公式}
\label{sec:vafa-intriligator-detail}

\subsection*{历史}

Vafa 和 Intriligator 在 1996 年基于物理中的镜像对称（mirror symmetry）思想提出了计算某些余维一子流形交点数的公式。

在旗流形 $Fl_n$ 的情形，这个公式给出了 Gromov-Witten 不变量的显式表达式。

\subsection*{公式陈述}

对于旗流形 $Fl_n$，Vafa-Intriligator 公式断言：
$$c_{u,v}^w(q) = \frac{1}{|W|} \sum_{\sigma \in W} \frac{\sigma(\sigma_u) \sigma(\sigma_v) \prod_{\alpha \in R^+} \alpha}{\prod_{\alpha \in R^+} (1 - q^\alpha \sigma(e^{-\alpha}))}$$

其中 $W$ 是 Weyl 群，$R^+$ 是正根集合。

这个公式的美妙之处在于：它将抽象的曲线计数问题转化为具体的求和。

\section{附录：量子 Schubert 多项式的正交性证明}
\label{sec:orthogonality-proof}

定理：量子 Schubert 多项式满足正交性关系
$$\langle \widetilde{\mathfrak{S}}_u, \widetilde{\mathfrak{S}}_v \rangle_Q = \delta_{v, w_0 u}$$

\textbf{证明思路}：

步骤 1：归约到典范元素。

由量子 Cauchy 恒等式，
$$\langle \widetilde{\mathfrak{S}}_{w_0}(x,y), \widetilde{\mathfrak{S}}_{w_0}(x,z) \rangle_Q^{(x)} = C(y,z)$$
其中右边是"典范元素"。

步骤 2：施加除差算子。

对两边施加除差算子 $\partial_{vw_0}^{(y)} \partial_{ww_0}^{(z)}$：
$$\langle \widetilde{\mathfrak{S}}_v(x,y), \widetilde{\mathfrak{S}}_w(x,z) \rangle_Q^{(x)} = \sum_{u \in S_n} \partial_{vw_0}^{(y)} \mathfrak{S}_u(y) \cdot \partial_{ww_0}^{(z)} \mathfrak{S}_{w_0 u}(z)$$

步骤 3：取极限。

取 $y=z=0$，利用经典的正交性关系 $\eta(\partial_w \mathfrak{S}_u) = \delta_{w,u}$，可得所求结果。

详细证明见 \cite[ Section 5]{KM}。

\section{附录：量子化映射的进一步性质}
\label{sec:quantization-further}

量子化映射 $\Phi: P_n \to \overline{P}_n$ 定义为
$$\Phi(f)(x) = \sum_{w \in S^{(n)}} \partial_w^{(y)} f(y) \widetilde{\mathfrak{S}}_w(x,y) \big|_{\overline{P}_n}$$

\textbf{性质汇总}：

\begin{enumerate}
\item \textbf{线性性}：$\Phi(af + bg) = a\Phi(f) + b\Phi(g)$
\item \textbf{配对保持}：$\langle \Phi(f), \Phi(g) \rangle_Q = \langle f, g \rangle$
\item \textbf{保持初等对称函数}：$\Phi(e_i) = \widetilde{e}_i$
\item \textbf{不保持乘法}：$\Phi(fg) \neq \Phi(f) \Phi(g)$
\end{enumerate}

\textbf{为什么不保持乘法？} 这是量子效应的核心体现。在量子理论中，乘法结构被量子修正了——这不是映射的缺陷，而是数学本质的改变。

\endinput
